% Generated draft: requires mathematical and source review.
\documentclass[11pt]{article}
\usepackage[margin=25mm]{geometry}
\usepackage{fontspec}
\setmainfont{Latin Modern Roman}
\usepackage{amsmath,amssymb,mathtools,unicode-math}
\setmathfont{Latin Modern Math}
\usepackage{enumitem,needspace,xurl,hyperref}
\setlength{\parindent}{0pt}
\setlength{\parskip}{5pt}
\setlength{\emergencystretch}{3em}
\newcommand{\N}{\mathbb{N}}
\newcommand{\Z}{\mathbb{Z}}
\newcommand{\Q}{\mathbb{Q}}
\newcommand{\R}{\mathbb{R}}
\newcommand{\C}{\mathbb{C}}
\newcommand{\F}{\mathbb{F}}
\newcommand{\A}{\mathbb{A}}
\newcommand{\PP}{\mathbb{P}}
\newcommand{\E}{\mathbb{E}}
\newcommand{\eps}{\varepsilon}
\newcommand{\dd}{\,\mathrm{d}}
\begin{document}
\section*{UnsolvedMath ID: 30000735}
\textbf{Source problem code:} catalogue entry
\subsection{Definitions and mathematical statement}
Consider the biased diameter games $D_d(a:b)$ on the edges of $K_n$ in the report. Under the usual Maker--Breaker interpretation, Maker and Breaker alternately claim previously unclaimed edges, $a$ and $b$ per turn respectively, and Maker wins if Maker's graph has diameter at most $d$. Graph distance in a disconnected graph is infinite. The source does not state turn order, end-of-board rules, or confirm the exact winning condition, so these conventions must be checked.

With the usual monotone bias convention, a breaking point may be denoted
\[
 b_{d,a}(n)=\min\{b\ge1:\text{Breaker has a winning strategy in }D_d(a:b)\}.
\]
The report's precise threshold convention must be recovered. The general conjectural request is that this breaking point be asymptotically close, up to a polylogarithmic factor, to the known Breaker-winning bias bounds for the parameter pairs considered in the report. Those bounds and the range of $(d,a)$ are not supplied, so no explicit general comparison function is invented.

The restored statement singles out two tasks:
\begin{enumerate}
\item Is the breaking point for $D_3(3:b)$ of order $n^{2/3}\operatorname{polylog}(n)$?
\item Determine the breaking point for $D_3(2:b)$.
\end{enumerate}
Here the unknown variable is Breaker's bias $b$ as $n\to\infty$, while $d=3$ and Maker's biases are fixed as stated. The source gives no candidate exponent for the second case.

The notation polylogarithmic is not precise enough to specify a single asymptotic equivalent: it may mean an unspecified power of $\log n$, or an allowed multiplicative loss above and below a reference bound. Its exponent, sign, constants, and the direction of the known comparison must be recovered. An assertion merely of the form $n^{2/3+o(1)}$ would not by itself preserve a fixed polylogarithmic-factor requirement. No exact logarithmic order or leading constant is silently chosen.
\subsection{Short English statement}
The question asks whether the bias threshold in the report's diameter games matches known Breaker-winning bounds up to logarithmic powers, especially the $n^{2/3}$ scale for Maker bias three. It separately asks for Maker bias two, without supplying a predicted scale.
\subsection{Sources}
\begin{itemize}
\item Local UnsolvedMath registry, permanent ID 30000735, code catalogue entry, ``Breaking Points in Diameter Positional Games.'' Proposed year 2007; proposer unspecified.
\item \emph{Mini-Workshop: Positional Games} (2008), Oberwolfach Reports, pp.\ 1073--1114, DOI 10.4171/owr/2007/20; \url{https://doi.org/10.4171/owr/2007/20}. The displayed report year 2008 differs from the DOI and proposed year 2007. This is the sole supplied literature reference.
\item Statement status corrected\_verified; the supplied verification says official-report content was checked and wording/notation normalized without changing the source-level claim. The retained original only states the general comparison; the explicit two cases occur in the restored field.
\item Literature triage checked 2026-08-21: open. It calls the break-point comparison conjectural and says no definitive resolution of this extracted claim was located. The exact game rules, bounds, and current status require independent verification. Created 2026-08-25; updated 2026-09-16.
\item Registry curation by Ulam AI / UnsolvedMath Contributors, CC BY 4.0; the underlying report retains its supplied provenance. Unreviewed statement-only draft.
\end{itemize}
\end{document}
